\section{Detection on the projective curve}\label{sec:detection}

Proposition~\ref{prop:integral} has constructed an integral class $c$
with $c|_{\Frac A}=0$ and coefficient image $i_*a_D$ on
\[
 T=\Spec A[1/\pi].
\]
By Lemma~\ref{lem:field}, $a_D=\beta_D\cdot[t]=\theta_D(t)$.
To prove its pushforward $i_*a_D$ nonzero, we extend the supported
class to the projective curve $C$, where proper pushforward permits
the norm test. The scheme $T$ is an inverse limit of
affine opens of $C$ with flat transition maps, and need not itself
be an open subset. The following lemma makes that passage explicit.

\begin{lemma}\label{lem:opens}
Let $\Sigma=R\setminus\mathfrak m$, where
$R=V[x,y]/(F)$ and $\mathfrak m=(\pi,x,y)$.  For $h\in\Sigma$ put
\[
 U_h=\Spec R[1/\pi,1/h]\subset C.
\]
Then the refinement maps $U_{hk}\to U_h$ are affine and flat, and
\begin{equation}\label{eq:continuity}
 T=\varprojlim_{h\in\Sigma}U_h,
 \qquad
 \varinjlim_{h\in\Sigma}G_3(U_h;\Z/5)
   \xrightarrow{\ \sim\ }G_3(T;\Z/5).
\end{equation}
Every $U_h$ contains $q_0$.  Its complement in $C$ is a finite set of
closed points $q$, and their residue fields $E_q=\kappa(q)$ satisfy
$5\mid[k_{E_q}:\F_5]$.
\end{lemma}

\begin{proof}
The multiplicative set $\Sigma$ is directed by taking products:
$hk$ gives a common refinement of $h$ and $k$.  Transitivity of
localization gives
\[
 \varinjlim_{h\in\Sigma}R[1/\pi,1/h]
 =R_{\mathfrak m}[1/\pi]=A[1/\pi].
\]
Each refinement is a principal localization, so the corresponding
scheme map is affine and flat.  The stages and their limit are affine
Noetherian schemes, hence separated.  Quillen's continuity theorem,
in its coherent $K'$ conclusion
\cite[\S7, Proposition~2.2]{QuillenHigher}, therefore applies to
their integral $G$-groups.

To pass to coefficients, use the natural short exact sequence
\[
 0\longrightarrow G_3(Y)/5G_3(Y)
 \longrightarrow G_3(Y;\Z/5)
 \longrightarrow G_2(Y)[5]\longrightarrow0
\]
for each stage $Y=U_h$ and for $Y=T$.  Filtered colimits of abelian
groups are exact, and thus commute with the quotient by five and
the kernel of multiplication by five.  Integral continuity in
degrees three and two identifies the outer terms in the resulting
comparison of short exact sequences.  It consequently identifies
the middle terms, proving the coefficient assertion
in~\eqref{eq:continuity}.

Since $h\notin\mathfrak m$, its image in the local divisor ring $S$
is a unit.  Thus $q_0=\Spec D$ belongs to $U_h$.  More generally,
consider a closed point $q\in C$ with residue field $E$ whose
extension $\Spec\OO_E\to\X$ specializes to $m$.
The inverse image of the chart $Z\ne0$ contains the closed point of
$\Spec\OO_E$, and hence is all of $\Spec\OO_E$.
On that chart the reduction of $h$ at $m$ is nonzero.  Its pullback
to $\OO_E$ is therefore a unit, so $q$ lies in $U_h$.
It follows that each point omitted from $U_h$ specializes away
from $m$.  Lemma~\ref{lem:model} then gives
$5\mid[k_E:\F_5]$.  Finally, $C$ is an integral projective curve
and $U_h$ is nonempty, so its closed complement consists of
finitely many closed points.
\end{proof}

\begin{proof}[Proof of Theorem~\ref{thm:main}]
Take the integral class $c\in K_3(A)$ constructed in
Proposition~\ref{prop:integral}, and put
\[
 w=(q_0)_*\theta_D(t)\in G_3(C;\Z/5).
\]
Coherent $G$-theory provides proper pushforward even when $C$ is singular.
The point immersion is finite, and the structural map
$p:C\to\Spec L$ is proper with $C$ carrying the ample line bundle
$\mathcal O_C(1)$.  Coherent pushforward and its composition law
therefore apply to both maps
\cite[\S7, paragraph 2.7 and the discussion following (2.8)]{QuillenHigher}.

Flat base change for the point immersion on each $U_h$
\cite[V.3.7.2]{Weibel}, followed by~\eqref{eq:continuity}, identifies
$w|_T$ with the pushforward of $\theta_D(t)$ along
$\Spec D\hookrightarrow T$.
By Proposition~\ref{prop:integral}, this is precisely the coefficient
reduction of $c|_T$.

Suppose that $w|_T=0$.  The element $w|_{U_1}$ represents this class
in the filtered colimit~\eqref{eq:continuity}.  A zero element in
that colimit becomes zero at a finite stage, so
$w|_{U_h}=0$ for some $h\in\Sigma$.
Give $Z_h=C\setminus U_h$ its reduced closed structure.
Apply Moore coefficients to coherent localization and d\'evissage
at these closed points \cite[\S7, Proposition~3.2 and
paragraph~3.4]{QuillenHigher}. The resulting exact segment is
\[
 \bigoplus_{q\in Z_h}K_3(E_q;\Z/5)
 \longrightarrow G_3(C;\Z/5)
 \longrightarrow G_3(U_h;\Z/5).
\]
Consequently there are classes $a_q\in K_3(E_q;\Z/5)$ such that
\begin{equation}\label{eq:omitted-relation}
 w=\sum_{q\in Z_h}q_*a_q.
\end{equation}
These are arbitrary coefficient classes supplied by localization.
Lemma~\ref{lem:field} applies to every finite extension $E_q/L$,
so its surjectivity writes each of them as
$a_q=\theta_{E_q}(u_q)$ for some $u_q\in E_q^\times$.

Apply the proper pushforward
$p_*:G_3(C;\Z/5)\to K_3(L;\Z/5)$ to
\eqref{eq:omitted-relation}.
For every closed point, the composite $p_*q_*$ is the field transfer.
Lemma~\ref{lem:opens} gives $5\mid[k_{E_q}:\F_5]$, and
Lemma~\ref{lem:norm-detector} therefore gives
\[
 \lambda_L\bigl(p_*q_*a_q\bigr)
 =v_L\bigl(N_{E_q/L}(u_q)\bigr)\bmod5
 =[k_{E_q}:\F_5]v_{E_q}(u_q)\bmod5=0.
\]
Thus \eqref{eq:omitted-relation} forces $\lambda_L(p_*w)=0$.
On the other hand, the distinguished point has residue field $D$,
and its parameter satisfies $N_{D/L}(t)=\pi$.  The same transfer
formula yields
\[
 \lambda_L(p_*w)
 =\lambda_L\bigl(\operatorname{Tr}_{D/L}\theta_D(t)\bigr)
 =v_L\bigl(N_{D/L}(t)\bigr)\bmod5
 =v_L(\pi)\bmod5=1.
\]
This contradiction proves $w|_T\ne0$.
Hence the coefficient reduction of $c|_T$ is nonzero, so the
integral class $c\in K_3(A)$ is itself nonzero.
Its image at the fraction field is zero by
Proposition~\ref{prop:integral}; together with the properties of $A$
in Lemma~\ref{lem:model}, this proves the theorem.
\end{proof}
